
\begingroup
\tikzset{external/figure name/.add={}{brewer_}}%
    \tikzset{
%        font=\scriptsize,
        matrix style/.style={
            font=\tiny,
            every matrix/.append style={
                matrix anchor=north,
                inner ysep=0pt,
%                fill=red!10,
                ampersand replacement=\&,
                matrix of nodes,
                nodes={
                    anchor=south,
                    minimum size=2mm,
                },
                row 1/.style={
                    inner sep=0pt,
                },
                column 1/.style={
                    inner ysep=0pt,
                    anchor=base east,
                },
                row 2 column 1/.style={
                    every node/.append style={
                        font=\bfseries\scriptsize,
%                        font=\bfseries\tiny,
                        draw,
                    },
                },
                row sep=0pt,
                column sep=0pt,
            },
        },
    }


% =============================================================================
%\def\MATRIXseq#1{%
\NewDocumentCommand{\MATRIXseq}{sm}{{%
\tikzset{/tikz/external/export=true}%
\newcommand*\scheme{#2}%
\pgfmanualpdflabel{colormap/\scheme-3}{}%
\pgfmanualpdflabel{colormap/\scheme-4}{}%
\pgfmanualpdflabel{colormap/\scheme-5}{}%
\pgfmanualpdflabel{colormap/\scheme-6}{}%
\pgfmanualpdflabel{colormap/\scheme-7}{}%
\pgfmanualpdflabel{colormap/\scheme-8}{}%
\pgfmanualpdflabel{colormap/\scheme-9}{}%
\pgfmanualpdflabel{colormap/\scheme}{}%
\foreach \x in {A,...,M} {\pgfmanualpdflabel{\scheme-\x}{}}%
\tikzsetnextfilename{brewer-seq-#2}%
\begin{tikzpicture}[
    matrix style,
]
    % instantiate all colormaps for \scheme:
    \pgfplotsset{
        colormap/\scheme-3,
        colormap/\scheme-4,
        colormap/\scheme-5,
        colormap/\scheme-6,
        colormap/\scheme-7,
        colormap/\scheme-8,
        colormap/\scheme-9,
        colormap/\scheme,
    }%
    %
    % directly access the correct color:
    % #1 is one of A,B,...,M
    \def\C##1{\node [fill=\scheme-##1] {};}%
    % in case the color should be white, also draw a black frame
    \def\D##1{\node [draw=black,fill=\scheme-##1] {};}%
    %
    % indirectly access the correct color
    % #1 is an index into the colormap, 0,1,...,13
    % #2 is the number of colors in the colormap (qualifies which colormap to use)
    % This is redundant, but it allows us to verify that the colormaps
    % work as intended without typos:
    \newcommand*\CC[2]{\node [/pgfplots/index of colormap=##1 of \scheme-##2, fill] {};}%
    % in case the color should be white, also draw a black frame
    \newcommand*\DD[2]{\node [/pgfplots/index of colormap=##1 of \scheme-##2, fill, draw=black] {};}%
    %
    % define a colorbar that has the same size as a row of colors in the matrix
    \newcommand*\CM{%
        \node[inner sep=0pt,xshift=-3.5pt,overlay,anchor=south west]{%
\IfBooleanTF{#1}{%
            \pgfmathparse{14*2mm + 16*\pgflinewidth + 1.0pt}%
}{%
            \pgfmathparse{14*2mm + 15*\pgflinewidth + 1.1pt}%
}
            \let\width=\pgfmathresult
            \pgfplotscolormaptoshadingspec{\scheme}{\width pt}\result
            \def\tempb{\pgfdeclarehorizontalshading{tempshading}{2mm}}%
            \expandafter\tempb\expandafter{\result}%
            \pgfuseshading{tempshading}%
        };
    }%
    \matrix {
\IfBooleanTF{#1}{%
                \& A         \& B         \& C         \& D         \& E         \& F         \& G         \& H         \& I         \& J         \& K         \& L         \& M         \\
        \scheme \& \D{A}     \& \C{B}     \& \C{C}     \& \C{D}     \& \C{E}     \& \C{F}     \& \C{G}     \& \C{H}     \& \C{I}     \& \C{J}     \& \C{K}     \& \C{L}     \& \C{M}     \\
        3       \&           \&           \& \CC{0}{3} \&           \&           \& \CC{1}{3} \&           \&           \& \CC{2}{3} \&           \&           \&           \&           \\
        4       \&           \& \CC{0}{4} \&           \&           \& \CC{1}{4} \&           \& \CC{2}{4} \&           \&           \& \CC{3}{4} \&           \&           \&           \\
        5       \&           \& \CC{0}{5} \&           \&           \& \CC{1}{5} \&           \& \CC{2}{5} \&           \& \CC{3}{5} \&           \& \CC{4}{5} \&           \&           \\
        6       \&           \& \CC{0}{6} \&           \& \CC{1}{6} \&           \& \CC{2}{6} \& \CC{3}{6} \&           \& \CC{4}{6} \&           \& \CC{5}{6} \&           \&           \\
        7       \&           \& \CC{0}{7} \&           \& \CC{1}{7} \&           \& \CC{2}{7} \& \CC{3}{7} \& \CC{4}{7} \&           \& \CC{5}{7} \&           \& \CC{6}{7} \&           \\
        8       \& \DD{0}{8} \&           \& \CC{1}{8} \& \CC{2}{8} \&           \& \CC{3}{8} \& \CC{4}{8} \& \CC{5}{8} \&           \& \CC{6}{8} \&           \& \CC{7}{8} \&           \\
        9       \& \DD{0}{9} \&           \& \CC{1}{9} \& \CC{2}{9} \&           \& \CC{3}{9} \& \CC{4}{9} \& \CC{5}{9} \&           \& \CC{6}{9} \& \CC{7}{9} \&           \& \CC{8}{9} \\
        CM      \& \CM       \&           \&           \&           \&           \&           \&           \&           \&           \&           \&           \&           \&           \\
%                \& A         \& B         \& C         \& D         \& E         \& F         \& G         \& H         \& I         \& J         \& K         \& L         \& M         \\
%                \& \D{A}     \& \C{B}     \& \C{C}     \& \C{D}     \& \C{E}     \& \C{F}     \& \C{G}     \& \C{H}     \& \C{I}     \& \C{J}     \& \C{K}     \& \C{L}     \& \C{M}     \\
%                \&           \&           \& \CC{0}{3} \&           \&           \& \CC{1}{3} \&           \&           \& \CC{2}{3} \&           \&           \&           \&           \\
%                \&           \& \CC{0}{4} \&           \&           \& \CC{1}{4} \&           \& \CC{2}{4} \&           \&           \& \CC{3}{4} \&           \&           \&           \\
%                \&           \& \CC{0}{5} \&           \&           \& \CC{1}{5} \&           \& \CC{2}{5} \&           \& \CC{3}{5} \&           \& \CC{4}{5} \&           \&           \\
%                \&           \& \CC{0}{6} \&           \& \CC{1}{6} \&           \& \CC{2}{6} \& \CC{3}{6} \&           \& \CC{4}{6} \&           \& \CC{5}{6} \&           \&           \\
%                \&           \& \CC{0}{7} \&           \& \CC{1}{7} \&           \& \CC{2}{7} \& \CC{3}{7} \& \CC{4}{7} \&           \& \CC{5}{7} \&           \& \CC{6}{7} \&           \\
%                \& \DD{0}{8} \&           \& \CC{1}{8} \& \CC{2}{8} \&           \& \CC{3}{8} \& \CC{4}{8} \& \CC{5}{8} \&           \& \CC{6}{8} \&           \& \CC{7}{8} \&           \\
%                \& \DD{0}{9} \&           \& \CC{1}{9} \& \CC{2}{9} \&           \& \CC{3}{9} \& \CC{4}{9} \& \CC{5}{9} \&           \& \CC{6}{9} \& \CC{7}{9} \&           \& \CC{8}{9} \\
}{%
                \& A         \& B         \& C         \& D         \& E         \& F         \& G         \& H         \& I         \& J         \& K         \& L         \& M         \\
        \scheme \& \C{A}     \& \C{B}     \& \C{C}     \& \C{D}     \& \C{E}     \& \C{F}     \& \C{G}     \& \C{H}     \& \C{I}     \& \C{J}     \& \C{K}     \& \C{L}     \& \C{M}     \\
        3       \&           \&           \& \CC{0}{3} \&           \&           \& \CC{1}{3} \&           \&           \& \CC{2}{3} \&           \&           \&           \&           \\
        4       \&           \& \CC{0}{4} \&           \&           \& \CC{1}{4} \&           \& \CC{2}{4} \&           \&           \& \CC{3}{4} \&           \&           \&           \\
        5       \&           \& \CC{0}{5} \&           \&           \& \CC{1}{5} \&           \& \CC{2}{5} \&           \& \CC{3}{5} \&           \& \CC{4}{5} \&           \&           \\
        6       \&           \& \CC{0}{6} \&           \& \CC{1}{6} \&           \& \CC{2}{6} \& \CC{3}{6} \&           \& \CC{4}{6} \&           \& \CC{5}{6} \&           \&           \\
        7       \&           \& \CC{0}{7} \&           \& \CC{1}{7} \&           \& \CC{2}{7} \& \CC{3}{7} \& \CC{4}{7} \&           \& \CC{5}{7} \&           \& \CC{6}{7} \&           \\
        8       \& \CC{0}{8} \&           \& \CC{1}{8} \& \CC{2}{8} \&           \& \CC{3}{8} \& \CC{4}{8} \& \CC{5}{8} \&           \& \CC{6}{8} \&           \& \CC{7}{8} \&           \\
        9       \& \CC{0}{9} \&           \& \CC{1}{9} \& \CC{2}{9} \&           \& \CC{3}{9} \& \CC{4}{9} \& \CC{5}{9} \&           \& \CC{6}{9} \& \CC{7}{9} \&           \& \CC{8}{9} \\
        CM      \& \CM       \&           \&           \&           \&           \&           \&           \&           \&           \&           \&           \&           \&           \\
%                \& A         \& B         \& C         \& D         \& E         \& F         \& G         \& H         \& I         \& J         \& K         \& L         \& M         \\
%                \& \C{A}     \& \C{B}     \& \C{C}     \& \C{D}     \& \C{E}     \& \C{F}     \& \C{G}     \& \C{H}     \& \C{I}     \& \C{J}     \& \C{K}     \& \C{L}     \& \C{M}     \\
%                \&           \&           \& \CC{0}{3} \&           \&           \& \CC{1}{3} \&           \&           \& \CC{2}{3} \&           \&           \&           \&           \\
%                \&           \& \CC{0}{4} \&           \&           \& \CC{1}{4} \&           \& \CC{2}{4} \&           \&           \& \CC{3}{4} \&           \&           \&           \\
%                \&           \& \CC{0}{5} \&           \&           \& \CC{1}{5} \&           \& \CC{2}{5} \&           \& \CC{3}{5} \&           \& \CC{4}{5} \&           \&           \\
%                \&           \& \CC{0}{6} \&           \& \CC{1}{6} \&           \& \CC{2}{6} \& \CC{3}{6} \&           \& \CC{4}{6} \&           \& \CC{5}{6} \&           \&           \\
%                \&           \& \CC{0}{7} \&           \& \CC{1}{7} \&           \& \CC{2}{7} \& \CC{3}{7} \& \CC{4}{7} \&           \& \CC{5}{7} \&           \& \CC{6}{7} \&           \\
%                \& \CC{0}{8} \&           \& \CC{1}{8} \& \CC{2}{8} \&           \& \CC{3}{8} \& \CC{4}{8} \& \CC{5}{8} \&           \& \CC{6}{8} \&           \& \CC{7}{8} \&           \\
%                \& \CC{0}{9} \&           \& \CC{1}{9} \& \CC{2}{9} \&           \& \CC{3}{9} \& \CC{4}{9} \& \CC{5}{9} \&           \& \CC{6}{9} \& \CC{7}{9} \&           \& \CC{8}{9} \\
}
    };
\end{tikzpicture}%
}}

\NewDocumentCommand{\MATRIXdiv}{sm}{{%
\newcommand*\scheme{#2}%
\tikzset{/tikz/external/export=true}%
\pgfmanualpdflabel{colormap/\scheme-3}{}%
\pgfmanualpdflabel{colormap/\scheme-4}{}%
\pgfmanualpdflabel{colormap/\scheme-5}{}%
\pgfmanualpdflabel{colormap/\scheme-6}{}%
\pgfmanualpdflabel{colormap/\scheme-7}{}%
\pgfmanualpdflabel{colormap/\scheme-8}{}%
\pgfmanualpdflabel{colormap/\scheme-9}{}%
\pgfmanualpdflabel{colormap/\scheme-10}{}%
\pgfmanualpdflabel{colormap/\scheme-11}{}%
\pgfmanualpdflabel{colormap/\scheme}{}%
\foreach \x in {A,...,O} {\pgfmanualpdflabel{\scheme-\x}{}}%
\tikzsetnextfilename{brewer-div-#2}%
\begin{tikzpicture}[
    matrix style,
]
    % instantiate all colormaps for \scheme:
    \pgfplotsset{
        colormap/\scheme-3,
        colormap/\scheme-4,
        colormap/\scheme-5,
        colormap/\scheme-6,
        colormap/\scheme-7,
        colormap/\scheme-8,
        colormap/\scheme-9,
        colormap/\scheme-10,
        colormap/\scheme-11,
        colormap/\scheme,
    }%
    %
    % directly access the correct color:
    % #1 is one of A,B,...,M
    \def\C##1{\node [fill=\scheme-##1] {};}%
    % in case the color should be white, also draw a black frame
    \def\D##1{\node [draw=black,fill=\scheme-##1] {};}%
    %
    % indirectly access the correct color
    % #1 is an index into the colormap, 0,1,...,13
    % #2 is the number of colors in the colormap (qualifies which colormap to use)
    % This is redundant, but it allows us to verify that the colormaps
    % work as intended without typos:
    \newcommand*\CC[2]{\node [/pgfplots/index of colormap=##1 of \scheme-##2, fill] {};}%
    % in case the color should be white, also draw a black frame
    \newcommand*\DD[2]{\node [/pgfplots/index of colormap=##1 of \scheme-##2, fill, draw=black] {};}%
    %
    % define a colorbar that has the same size as a row of colors in the matrix
    \newcommand*\CM{%
        \node[inner sep=0pt,xshift=-3.5pt,overlay,anchor=south west]{%
\IfBooleanTF{#1}{%
            \pgfmathparse{16*2mm + 18*\pgflinewidth + 2.3pt}%
}{%
            \pgfmathparse{16*2mm + 17*\pgflinewidth + 2.3pt}%
}%
            \let\width=\pgfmathresult
            \pgfplotscolormaptoshadingspec{\scheme}{\width pt}\result
            \def\tempb{\pgfdeclarehorizontalshading{tempshading}{2mm}}%
            \expandafter\tempb\expandafter{\result}%
            \pgfuseshading{tempshading}%
        };
    }%
    \matrix {
\IfBooleanTF{#1}{%
                \& A          \& B          \& C         \& D          \& E         \& F          \& G          \& H          \& I          \& J          \& K         \& L          \& M         \& N          \& O           \\
        \scheme \& \C{A}      \& \C{B}      \& \C{C}     \& \C{D}      \& \C{E}     \& \C{F}      \& \C{G}      \& \D{H}      \& \C{I}      \& \C{J}      \& \C{K}     \& \C{L}      \& \C{M}     \& \C{N}      \& \C{O}       \\
        3       \&            \&            \&           \&            \& \CC{0}{3} \&            \&            \& \DD{1}{3}  \&            \&            \& \CC{2}{3} \&            \&           \&            \&             \\
        4       \&            \&            \& \CC{0}{4} \&            \&           \& \CC{1}{4}  \&            \&            \&            \& \CC{2}{4}  \&           \&            \& \CC{3}{4} \&            \&             \\
        5       \&            \&            \& \CC{0}{5} \&            \&           \& \CC{1}{5}  \&            \& \DD{2}{5}  \&            \& \CC{3}{5}  \&           \&            \& \CC{4}{5} \&            \&             \\
        6       \&            \& \CC{0}{6}  \&           \&            \& \CC{1}{7} \&            \& \CC{2}{6}  \&            \& \CC{3}{6}  \&            \& \CC{4}{6} \&            \&           \& \CC{5}{6}  \&             \\
        7       \&            \& \CC{0}{7}  \&           \&            \& \CC{1}{7} \&            \& \CC{2}{7}  \& \DD{3}{7}  \& \CC{4}{7}  \&            \& \CC{5}{7} \&            \&           \& \CC{6}{7}  \&             \\
        8       \&            \& \CC{0}{8}  \&           \& \CC{1}{8}  \&           \& \CC{2}{8}  \& \CC{3}{8}  \&            \& \CC{4}{8}  \& \CC{5}{8}  \&           \& \CC{6}{8}  \&           \& \CC{7}{8}  \&             \\
        9       \&            \& \CC{0}{9}  \&           \& \CC{1}{9}  \&           \& \CC{2}{9}  \& \CC{3}{9}  \& \DD{4}{9}  \& \CC{5}{9}  \& \CC{6}{9}  \&           \& \CC{7}{9}  \&           \& \CC{8}{9}  \&             \\
        10      \& \CC{0}{10} \& \CC{1}{10} \&           \& \CC{2}{10} \&           \& \CC{3}{10} \& \CC{4}{10} \&            \& \CC{5}{10} \& \CC{6}{10} \&           \& \CC{7}{10} \&           \& \CC{8}{10} \& \CC{9}{10}  \\
        11      \& \CC{0}{11} \& \CC{1}{11} \&           \& \CC{2}{11} \&           \& \CC{3}{11} \& \CC{4}{11} \& \DD{5}{11} \& \CC{6}{11} \& \CC{7}{11} \&           \& \CC{8}{11} \&           \& \CC{9}{11} \& \CC{10}{11} \\
        CM      \& \CM        \&            \&           \&            \&           \&            \&            \&            \&            \&            \&           \&            \&           \&            \&             \\
}{%
                \& A          \& B          \& C         \& D          \& E         \& F          \& G          \& H          \& I          \& J          \& K         \& L          \& M         \& N          \& O           \\
        \scheme \& \C{A}      \& \C{B}      \& \C{C}     \& \C{D}      \& \C{E}     \& \C{F}      \& \C{G}      \& \C{H}      \& \C{I}      \& \C{J}      \& \C{K}     \& \C{L}      \& \C{M}     \& \C{N}      \& \C{O}       \\
        3       \&            \&            \&           \&            \& \CC{0}{3} \&            \&            \& \CC{1}{3}  \&            \&            \& \CC{2}{3} \&            \&           \&            \&             \\
        4       \&            \&            \& \CC{0}{4} \&            \&           \& \CC{1}{4}  \&            \&            \&            \& \CC{2}{4}  \&           \&            \& \CC{3}{4} \&            \&             \\
        5       \&            \&            \& \CC{0}{5} \&            \&           \& \CC{1}{5}  \&            \& \CC{2}{5}  \&            \& \CC{3}{5}  \&           \&            \& \CC{4}{5} \&            \&             \\
        6       \&            \& \CC{0}{6}  \&           \&            \& \CC{1}{7} \&            \& \CC{2}{6}  \&            \& \CC{3}{6}  \&            \& \CC{4}{6} \&            \&           \& \CC{5}{6}  \&             \\
        7       \&            \& \CC{0}{7}  \&           \&            \& \CC{1}{7} \&            \& \CC{2}{7}  \& \CC{3}{7}  \& \CC{4}{7}  \&            \& \CC{5}{7} \&            \&           \& \CC{6}{7}  \&             \\
        8       \&            \& \CC{0}{8}  \&           \& \CC{1}{8}  \&           \& \CC{2}{8}  \& \CC{3}{8}  \&            \& \CC{4}{8}  \& \CC{5}{8}  \&           \& \CC{6}{8}  \&           \& \CC{7}{8}  \&             \\
        9       \&            \& \CC{0}{9}  \&           \& \CC{1}{9}  \&           \& \CC{2}{9}  \& \CC{3}{9}  \& \CC{4}{9}  \& \CC{5}{9}  \& \CC{6}{9}  \&           \& \CC{7}{9}  \&           \& \CC{8}{9}  \&             \\
        10      \& \CC{0}{10} \& \CC{1}{10} \&           \& \CC{2}{10} \&           \& \CC{3}{10} \& \CC{4}{10} \&            \& \CC{5}{10} \& \CC{6}{10} \&           \& \CC{7}{10} \&           \& \CC{8}{10} \& \CC{9}{10}  \\
        11      \& \CC{0}{11} \& \CC{1}{11} \&           \& \CC{2}{11} \&           \& \CC{3}{11} \& \CC{4}{11} \& \CC{5}{11} \& \CC{6}{11} \& \CC{7}{11} \&           \& \CC{8}{11} \&           \& \CC{9}{11} \& \CC{10}{11} \\
        CM      \& \CM        \&            \&           \&            \&           \&            \&            \&            \&            \&            \&           \&            \&           \&            \&             \\
}
    };
\end{tikzpicture}%
}}

\NewDocumentCommand{\MATRIXqual}{mmO{0}}{{%
\tikzset{/tikz/external/export=true}%
\newcommand*\scheme{#2}%
\pgfmanualpdflabel{colormap/\scheme-3}{}%
\pgfmanualpdflabel{colormap/\scheme-4}{}%
\pgfmanualpdflabel{colormap/\scheme-5}{}%
\pgfmanualpdflabel{colormap/\scheme-6}{}%
\pgfmanualpdflabel{colormap/\scheme-7}{}%
\pgfmanualpdflabel{colormap/\scheme-8}{}%
\ifnum#1>8
\pgfmanualpdflabel{colormap/\scheme-9}{}%
\ifnum#1>9
\pgfmanualpdflabel{colormap/\scheme-10}{}%
\pgfmanualpdflabel{colormap/\scheme-11}{}%
\pgfmanualpdflabel{colormap/\scheme-12}{}%
\fi
\fi
\pgfmanualpdflabel{colormap/\scheme}{}%
\tikzsetnextfilename{brewer-qualit-#2}%
\begin{tikzpicture}[
    baseline,
    matrix style,
]
    \newcommand*\classes{#1}%
    % instantiate all colormaps for \scheme:
    \pgfplotsset{
        colormap/\scheme-3,
        colormap/\scheme-4,
        colormap/\scheme-5,
        colormap/\scheme-6,
        colormap/\scheme-7,
        colormap/\scheme-8,
    }
\foreach \x in {A,...,H} {\pgfmanualpdflabel{\scheme-\x}{}}%
\ifnum\classes>8
\foreach \x in {I} {\pgfmanualpdflabel{\scheme-\x}{}}%
    \pgfplotsset{
        colormap/\scheme-9,
    }
\fi
\ifnum\classes=12
\foreach \x in {J,...,L} {\pgfmanualpdflabel{\scheme-\x}{}}%
    \pgfplotsset{
        colormap/\scheme-10,
        colormap/\scheme-11,
        colormap/\scheme-12,
    }
\fi
    \pgfplotsset{
        colormap/\scheme,
    }%
    %
    \newcommand*\EmptyRows{#3}

    % directly access the correct color:
    % #1 is one of A,B,...,M
    \def\C##1{\node [fill=\scheme-##1] {};}
    %
    % indirectly access the correct color
    % #1 is an index into the colormap, 0,1,...,13
    % #2 is the number of colors in the colormap (qualifies which colormap to use)
    % This is redundant, but it allows us to verify that the colormaps
    % work as intended without typos:
    \newcommand*\CC[2]{\node [/pgfplots/index of colormap=##1 of \scheme-##2, fill] {};}
    %
    % define dummy node to just produce space
    \newcommand*\Dummy{\node {};}
%    % define a colorbar that has the same size as a row of colors in the matrix
%\ifnum\classes=8
%    \newcommand*\CM{%
%        \node[inner sep=0pt,xshift=-3.5pt,overlay,anchor=south west]{%
%                \pgfmathparse{8*2mm + 8pt}%
%            \let\width=\pgfmathresult%
%            \pgfplotscolormaptoshadingspec{\scheme}{\width pt}\result%
%            \def\tempb{\pgfdeclarehorizontalshading{tempshading}{2mm}}%
%            \expandafter\tempb\expandafter{\result}%
%            \pgfuseshading{tempshading}%
%        };
%    }
%\else\ifnum\classes=9
%    \newcommand*\CM{%
%        \node[inner sep=0pt,xshift=-3.5pt,overlay,anchor=south west]{%
%                \pgfmathparse{9*2mm + 10*\pgflinewidth + 2.3pt}%
%            \let\width=\pgfmathresult%
%            \pgfplotscolormaptoshadingspec{\scheme}{\width pt}\result%
%            \def\tempb{\pgfdeclarehorizontalshading{tempshading}{2mm}}%
%            \expandafter\tempb\expandafter{\result}%
%            \pgfuseshading{tempshading}%
%        };
%    }
%\else\ifnum\classes=12
%    \newcommand*\CM{%
%        \node[inner sep=0pt,xshift=-3.5pt,overlay,anchor=south west]{%
%                \pgfmathparse{12*2mm + 13*\pgflinewidth + 2.3pt}%
%            \let\width=\pgfmathresult%
%            \pgfplotscolormaptoshadingspec{\scheme}{\width pt}\result%
%            \def\tempb{\pgfdeclarehorizontalshading{tempshading}{2mm}}%
%            \expandafter\tempb\expandafter{\result}%
%            \pgfuseshading{tempshading}%
%        };
%    }
%\fi\fi\fi
\ifnum\classes=8
    \matrix (m) {
                \& A          \& B          \& C          \& D          \& E          \& F          \& G          \& H          \\
        \scheme \& \C{A}      \& \C{B}      \& \C{C}      \& \C{D}      \& \C{E}      \& \C{F}      \& \C{G}      \& \C{H}      \\
        3       \& \CC{0}{3}  \& \CC{1}{3}  \& \CC{2}{3}  \&            \&            \&            \&            \&            \\
        4       \& \CC{0}{4}  \& \CC{1}{4}  \& \CC{2}{4}  \& \CC{3}{4}  \&            \&            \&            \&            \\
        5       \& \CC{0}{5}  \& \CC{1}{5}  \& \CC{2}{5}  \& \CC{3}{5}  \& \CC{4}{5}  \&            \&            \&            \\
        6       \& \CC{0}{6}  \& \CC{1}{6}  \& \CC{2}{6}  \& \CC{3}{6}  \& \CC{4}{6}  \& \CC{5}{6}  \&            \&            \\
        7       \& \CC{0}{7}  \& \CC{1}{7}  \& \CC{2}{7}  \& \CC{3}{7}  \& \CC{4}{7}  \& \CC{5}{7}  \& \CC{6}{7}  \&            \\
        8       \& \CC{0}{8}  \& \CC{1}{8}  \& \CC{2}{8}  \& \CC{3}{8}  \& \CC{4}{8}  \& \CC{5}{8}  \& \CC{6}{8}  \& \CC{7}{8}  \& \Dummy \& \Dummy \& \Dummy \& \Dummy \\
%%%%%        CM      \& \CM        \&            \&            \&            \&            \&            \&            \&            \\
    };
\else\ifnum\classes=9
    \matrix (m) {
                \& A          \& B          \& C          \& D          \& E          \& F          \& G          \& H          \& I         \\
        \scheme \& \C{A}      \& \C{B}      \& \C{C}      \& \C{D}      \& \C{E}      \& \C{F}      \& \C{G}      \& \C{H}      \& \C{I}     \\
        3       \& \CC{0}{3}  \& \CC{1}{3}  \& \CC{2}{3}  \&            \&            \&            \&            \&            \&           \\
        4       \& \CC{0}{4}  \& \CC{1}{4}  \& \CC{2}{4}  \& \CC{3}{4}  \&            \&            \&            \&            \&           \\
        5       \& \CC{0}{5}  \& \CC{1}{5}  \& \CC{2}{5}  \& \CC{3}{5}  \& \CC{4}{5}  \&            \&            \&            \&           \\
        6       \& \CC{0}{6}  \& \CC{1}{6}  \& \CC{2}{6}  \& \CC{3}{6}  \& \CC{4}{6}  \& \CC{5}{6}  \&            \&            \&           \\
        7       \& \CC{0}{7}  \& \CC{1}{7}  \& \CC{2}{7}  \& \CC{3}{7}  \& \CC{4}{7}  \& \CC{5}{7}  \& \CC{6}{7}  \&            \&           \\
        8       \& \CC{0}{8}  \& \CC{1}{8}  \& \CC{2}{8}  \& \CC{3}{8}  \& \CC{4}{8}  \& \CC{5}{8}  \& \CC{6}{8}  \& \CC{7}{8}  \&           \\
        9       \& \CC{0}{9}  \& \CC{1}{9}  \& \CC{2}{9}  \& \CC{3}{9}  \& \CC{4}{9}  \& \CC{5}{9}  \& \CC{6}{9}  \& \CC{7}{9}  \& \CC{8}{9} \& \Dummy \& \Dummy \& \Dummy \\
%%%%%        CM      \& \CM        \&            \&            \&            \&            \&            \&            \&            \&           \\
    };
\else\ifnum\classes=12
    \matrix (m) {
                \& A          \& B          \& C          \& D          \& E          \& F          \& G          \& H          \& I          \& J          \& K           \& L           \\
        \scheme \& \C{A}      \& \C{B}      \& \C{C}      \& \C{D}      \& \C{E}      \& \C{F}      \& \C{G}      \& \C{H}      \& \C{I}      \& \C{J}      \& \C{K}       \& \C{L}       \\
        3       \& \CC{0}{3}  \& \CC{1}{3}  \& \CC{2}{3}  \&            \&            \&            \&            \&            \&            \&            \&             \&             \\
        4       \& \CC{0}{4}  \& \CC{1}{4}  \& \CC{2}{4}  \& \CC{3}{4}  \&            \&            \&            \&            \&            \&            \&             \&             \\
        5       \& \CC{0}{5}  \& \CC{1}{5}  \& \CC{2}{5}  \& \CC{3}{5}  \& \CC{4}{5}  \&            \&            \&            \&            \&            \&             \&             \\
        6       \& \CC{0}{6}  \& \CC{1}{6}  \& \CC{2}{6}  \& \CC{3}{6}  \& \CC{4}{6}  \& \CC{5}{6}  \&            \&            \&            \&            \&             \&             \\
        7       \& \CC{0}{7}  \& \CC{1}{7}  \& \CC{2}{7}  \& \CC{3}{7}  \& \CC{4}{7}  \& \CC{5}{7}  \& \CC{6}{7}  \&            \&            \&            \&             \&             \\
        8       \& \CC{0}{8}  \& \CC{1}{8}  \& \CC{2}{8}  \& \CC{3}{8}  \& \CC{4}{8}  \& \CC{5}{8}  \& \CC{6}{8}  \& \CC{7}{8}  \&            \&            \&             \&             \\
        9       \& \CC{0}{9}  \& \CC{1}{9}  \& \CC{2}{9}  \& \CC{3}{9}  \& \CC{4}{9}  \& \CC{5}{9}  \& \CC{6}{9}  \& \CC{7}{9}  \& \CC{8}{9}  \&            \&             \&             \\
        10      \& \CC{0}{10} \& \CC{1}{10} \& \CC{2}{10} \& \CC{3}{10} \& \CC{4}{10} \& \CC{5}{10} \& \CC{6}{10} \& \CC{7}{10} \& \CC{8}{10} \& \CC{9}{10} \&             \&             \\
        11      \& \CC{0}{11} \& \CC{1}{11} \& \CC{2}{11} \& \CC{3}{11} \& \CC{4}{11} \& \CC{5}{11} \& \CC{6}{11} \& \CC{7}{11} \& \CC{8}{11} \& \CC{9}{11} \& \CC{10}{11} \&             \\
        12      \& \CC{0}{12} \& \CC{1}{12} \& \CC{2}{12} \& \CC{3}{12} \& \CC{4}{12} \& \CC{5}{12} \& \CC{6}{12} \& \CC{7}{12} \& \CC{8}{12} \& \CC{9}{12} \& \CC{10}{12} \& \CC{11}{12} \\
%%%%%        CM      \& \CM        \&            \&            \&            \&            \&            \&            \&            \&            \&            \&             \&             \\
    };
\fi\fi\fi

% insert space as fake empty rows
\ifnum\EmptyRows>0
    \path node [below=\EmptyRows*2mm of m,anchor=south] {};
\fi

\end{tikzpicture}%
}}

\section{ColorBrewer}
\def\pgfplotsmanualcurlibrary{colormaps}

{\emph{由 Vincent A.\ Traag 与 Stefan Pinnow 扩展}}

\begin{pgfplotslibrary}{colorbrewer}

本库将 Cynthia Brewer 创建、发布于
\url{http://colorbrewer2.org/} 的配色方案引入 PGFPlots。这些方案最初是为制图需求
设计的，但如今也用于其他类型的数据可视化。ColorBrewer 方案分为以下类型
%
\begin{description}
    \item[sequential] 适用于从低到高有序的数据；
    \item[diverging] 用于突出相对于某个均值的变化；
    \item[qualitative] 其中颜色没有特殊顺序。
\end{description}
%
每个方案最多可包含 12 个不同的数据类（即颜色），并同时以色图与循环列表的形式
提供。

{%
\centering

\pgfplotsset{
    brewer example/.style={
        small,
        title style={font=\normalfont},
    }
}%
\renewcommand{\arraystretch}{2}%
\tikzset{/tikz/external/export=true}%
\begin{tabular}{cc}
\begin{tikzpicture}[baseline]
    \begin{axis}[brewer example, hide axis,title={Diverging \texttt{colormap/PiYG-11}},colormap/PiYG-11,colorbar left]
        \addplot3[surf,
            domain=0:360,samples=40]
                {sin(x)*sin(y)};
    \end{axis}
\end{tikzpicture}
&
% ----- example in section 4.3 on page 43 -----
\tikzset{/tikz/external/export=true}%
\begin{tikzpicture}[baseline]
    % --- changed by Mo-Gul ---
    \begin{axis}[
        brewer example,
        hide axis,
        view={60}{30},
        title={Sequential \texttt{colormap/PuBu-9}},
        colormap/PuBu-9,
        colorbar right,
        clip bounding box=default tikz,
    ]
    % -------------------------
        \addplot3[
            surf,shader=flat,
            samples=20,
            domain=-1:0,y domain=0:2*pi,
            z buffer=sort
        ] (
            { sqrt(1-x^2) * cos(deg(y)) },
            { sqrt(1-x^2) * sin(deg(y)) },
            x
        );
    \end{axis}
\end{tikzpicture}%
% ---------------------------------------------
\\
% ----- second last example in section 4.6.9 on page 159 -----
% Preamble: \pgfplotsset{width=7cm,compat=1.13}
\tikzset{/tikz/external/export=true}%
\begin{tikzpicture}[baseline]
    % --- changed by Mo-Gul ---
    \begin{axis}[
        brewer example,
        hide axis,
        view={60}{30},
        title={Diverging \texttt{colormap/RdGy-11}},
        colormap/RdGy-11,
        colorbar left,
        clip bounding box=default tikz,
    ]
    % -------------------------
        \addplot3 [
            mesh,z buffer=sort,
            scatter,only marks,scatter src=z,
            samples=30,domain=-1:1,y domain=0:2*pi,
        ] (
            { sqrt(1-x^2) * cos(deg(y)) },
            { sqrt(1-x^2) * sin(deg(y)) },
            x
        );
    \end{axis}
\end{tikzpicture}
&
\begin{tikzpicture}[baseline]
    \begin{axis}[
        brewer example,
        hide axis,
        title={Sequential \texttt{colormap/YlOrBr}},
        colormap/YlOrBr,
        colorbar right,
        clip bounding box=default tikz,
    ]
    \addplot3[surf,samples=9,domain=0:1]
        {(1-abs(2*(x-0.5))) * (1-abs(2*(y-0.5)))};
    %\addlegendentry{$\phi_x \phi_y$}
    %\addplot3+[ultra thick] coordinates {(0,0,0) (0.5,0,1) (1,0,0)};
    %\addlegendentry{$\phi_x $}
    %\addplot3+[ultra thick] coordinates {(1,0,0) (1,0.5,1) (1,1,0)};
    %\addlegendentry{$\phi_y $}
    \end{axis}
\end{tikzpicture}
\\
%
% ------------------------------------------------------------
%
% ----- first example in section 4.7.4 on page 180 -----
% Preamble: \pgfplotsset{width=7cm,compat=1.13}
\begin{tikzpicture}[baseline]
\begin{loglogaxis}[
    brewer example,
    legend pos=south west,
    xlabel=\textsc{Dof},
    ylabel=$L_2$ Error,
    title={Qualitative \texttt{cycle list/Set1-5}},
    cycle list/Set1-5,
    hide axis,
    thick,
    legend style={draw=none},
    cycle multiindex* list={
        mark list*\nextlist
        Set1-5\nextlist
    },
        clip bounding box=default tikz,
    % -----------------------
]
\addplot coordinates {
    (5,8.312e-02)    (17,2.547e-02)   (49,7.407e-03)
    (129,2.102e-03)  (321,5.874e-04)  (769,1.623e-04)
    (1793,4.442e-05) (4097,1.207e-05) (9217,3.261e-06)
};
\addplot coordinates {
    (7,8.472e-02)    (31,3.044e-02)   (111,1.022e-02)
    (351,3.303e-03)  (1023,1.039e-03) (2815,3.196e-04)
    (7423,9.658e-05) (18943,2.873e-05)
    (47103,8.437e-06)
};
\addplot coordinates {
    (9,7.881e-02)     (49,3.243e-02)  (209,1.232e-02)
    (769,4.454e-03)   (2561,1.551e-03)
    (7937,5.236e-04)  (23297,1.723e-04)
    (65537,5.545e-05) (178177,1.751e-05)};
\addplot coordinates {
    (11,6.887e-02)    (71,3.177e-02) (351,1.341e-02)
    (1471,5.334e-03)  (5503,2.027e-03)
    (18943,7.415e-04) (61183,2.628e-04)
    (187903,9.063e-05) (553983,3.053e-05)};

\addplot coordinates {
    (13,5.755e-02)    (97,2.925e-02) (545,1.351e-02)
    (2561,5.842e-03)  (10625,2.397e-03)
    (40193,9.414e-04) (141569,3.564e-04)
    (471041,1.308e-04) (1496065,4.670e-05)
};
\legend{$d=2$,$d=3$,$d=4$,$d=5$,$d=6$}
\end{loglogaxis}
\end{tikzpicture}
&
% ----- example in section 4.7.7 on page 195 -----
\begin{tikzpicture}[baseline]
\pgfplotsset{
    cycle from colormap manual style/.style={
        x=3cm,y=10pt,ytick=\empty,
        colorbar style={x=,y=,ytick=\empty},
        point meta min=0,point meta max=1,
        stack plots=y,
        y dir=reverse,colorbar style={y dir=reverse},
        every axis plot/.style={line width=2pt},
        legend entries={0,...,20},
        legend pos=outer north east,
    }
}
\begin{axis}[
    % --- changed by Mo-Gul -----
    % also the text before and after that code has to be adjusted
    brewer example,
    hide axis,
    title={Diverging \texttt{cycle list/RdYlBu-4}},
    cycle list/RdYlBu-4,
    %
    cycle from colormap manual style,
    x=,y=,
    legend entries=,
        clip bounding box=default tikz,
]
    \addplot coordinates {(0,1) (0.5,1) (1,1)};
    \addplot coordinates {(0,1) (0.5,1) (1,1)};
    \addplot coordinates {(0,1) (0.5,1) (1,1)};
    \addplot coordinates {(0,1) (0.5,1) (1,1)};
    \addplot coordinates {(0,1) (0.5,1) (1,1)};
    \addplot coordinates {(0,1) (0.5,1) (1,1)};
    \addplot coordinates {(0,1) (0.5,1) (1,1)};
    \addplot coordinates {(0,1) (0.5,1) (1,1)};
    \addplot coordinates {(0,1) (0.5,1) (1,1)};
    \addplot coordinates {(0,1) (0.5,1) (1,1)};
    \addplot coordinates {(0,1) (0.5,1) (1,1)};
\end{axis}
\end{tikzpicture}\\
\end{tabular}

}%
% ------------------------------------------------


\subsection{用法}

下面以“色样”的图形形式展示可用的配色方案。


\subsubsection*{配色方案“色样”}
\label{sec:pgfplots:brewer:usage}

色样是一个矩阵，展示某一特定方案的所有可用颜色以及可用的颜色组合。

{\centering
% inspired from http://mkweb.bcgsc.ca/brewer/swatches/brewer-palettes-swatches.pdf
\tikzset{/tikz/external/export=true}%
\tikzsetnextfilename{brewer-fig-intro-GnBu}%
\begin{tikzpicture}[
    font=\scriptsize,
]
    \newcommand*\scheme{GnBu}
    % instantiate all colormaps for \scheme:
    \pgfplotsset{
        colormap/\scheme-3,
        colormap/\scheme-4,
        colormap/\scheme-5,
        colormap/\scheme-6,
        colormap/\scheme-7,
        colormap/\scheme-8,
        colormap/\scheme-9,
        colormap/\scheme,
    }%

    % directly access the correct color:
    % #1 is one of A,B,...,M
    \def\C#1{\node [fill=\scheme-#1] {};}

    % just set a label to an invisible cell
    \def\D#1{\node [draw=none] (#1) {};}

    % define a colorbar that has the same size as a row of colors in the matrix
    \newcommand*\CM{%
        \node[inner sep=0pt,xshift=-4.4pt,overlay,anchor=south west]{%
            \pgfmathparse{13*3mm + 0pt}%
            \let\width=\pgfmathresult
            \pgfplotscolormaptoshadingspec{\scheme}{\width pt}\result
            \def\tempb{\pgfdeclarehorizontalshading{tempshading}{3mm}}%
            \expandafter\tempb\expandafter{\result}%
            \pgfuseshading{tempshading}%
        };
    }

    \matrix[
        ampersand replacement=\&,
        matrix of nodes,
        nodes={
            anchor=south,
            minimum size=3mm,
        },
        row 1/.style={
            inner sep=0pt,
        },
        column 1/.style={
            inner ysep=0pt,
            anchor=base east,
        },
        row 2 column 1/.style={
            every node/.append style={
                font=\bfseries\scriptsize,
                draw,
            },
        },
        row sep=0pt,
        column sep=0pt,
    ] (m) {
                \& A     \& B     \& C     \& D     \& E     \& F     \& G     \& H     \& I     \& J     \& K     \& L     \& M     \\
        \scheme \& \C{A} \& \C{B} \& \C{C} \& \C{D} \& \C{E} \& \C{F} \& \C{G} \& \C{H} \& \C{I} \& \C{J} \& \C{K} \& \C{L} \& \node [fill=\scheme-M] (c) {}; \\
        3       \&       \&       \& \C{C} \&       \&       \& \C{F} \&       \&       \& \C{I} \&       \&       \&       \&       \\
        4       \&       \& \C{B} \&       \&       \& \C{E} \&       \& \C{G} \&       \&       \& \C{J} \&       \&       \&       \\
        5       \&       \& \C{B} \&       \&       \& \C{E} \&       \& \C{G} \&       \& \C{I} \&       \& \C{K} \&       \&       \\
        6       \&       \& \C{B} \&       \& \C{D} \&       \& \C{F} \& \C{G} \&       \& \C{I} \&       \& \C{K} \&       \&       \\
        7       \& \D{a} \& \C{B} \&       \& \C{D} \&       \& \C{F} \& \C{G} \& \C{H} \&       \& \C{J} \&       \& \C{L} \& \D{b} \\
        8       \& \C{A} \&       \& \C{C} \& \C{D} \&       \& \C{F} \& \C{G} \& \C{H} \&       \& \C{J} \&       \& \C{L} \&       \\
        9       \& \C{A} \&       \& \C{C} \& \C{D} \&       \& \C{F} \& \C{G} \& \C{H} \&       \& \C{J} \& \C{K} \&       \& \C{M} \\
        CM      \& \CM   \&       \&       \&       \&       \&       \&       \&       \&       \&       \&       \&       \& \D{CM} \\
    };

    \node [coordinate,pin={[align=center]above left:(short) \\ scheme \\ name}]
        at ([xshift=5pt] m-2-1.north west) {};
    \node [coordinate,pin={right:letters to create scheme color name}]
        at ([xshift=5pt] m-1-14) {};
    \node [coordinate,pin={right:all colors of scheme}]
        at ([xshift=5pt] c) {};
    \node [coordinate,pin={[align=center]below:numbers to create \\ (full) scheme name \\ (number of data classes)}]
        at ([xshift=0pt] m-10-1.south) {};
    \node [coordinate,pin={[align=left]right:colormap based on \\ previous row colors; \\
                                             (also) accessible by the \\ short scheme name}]
        at ([xshift=5pt] CM) {};

    \draw [red,rounded corners=2pt]
        ([xshift=-1pt] a.north west)
            rectangle
        ([xshift=+1pt] b.south east)
    ;
    \node [coordinate,pin=right:colors of scheme \texttt{GnBu-7}]
        at ([xshift=1pt] b.east) {};

\end{tikzpicture}

}%

此类色样的阅读方式如下：
%
\begin{enumerate}
    \item 在该色块的左上角，可以找到方案的（短）名称。
    \item 随后的色块是该方案所包含的颜色。
    \item 要得到（完整）颜色名，将（短）方案名与第一行中的一个字母用连字符连接，
        例如

        |\tikz \fill[color=GnBu-H] (0,0) rectangle (1em,2ex);|

        其结果为 \tikz \fill[color=GnBu-H] (0,0) rectangle
        (1em,2ex);。
    \item 要得到完整方案名（可用作色图或循环列表），将（短）方案名与第一列中的
        一个数字用连字符连接，例如 |cycle list name=GnBu-7|。\\
        也提供了仅使用（短）方案名的特殊情况，它是数字最大的完整方案名的别名，
        例如方案名 \texttt{GnBu} 等价于 \texttt{GnBu-9}。
    \item 最后一行展示基于前一行颜色的连续色图，在本例中即 \texttt{GnBu-9}。
    \item 矩阵的其余部分展示对应方案中使用的颜色。
\end{enumerate}


\subsubsection{启用配色方案}

要启用某个 |colorbrewer| |colormap|，例如 |BuGn-5|，需使用键

    |colormap/BuGn-5|。

\noindent 这会初始化并选中相应的 |colormap|。它也会初始化相应的 |cycle list|
（但不会选中）。要初始化并选中名为 |BuGn-5| 的 |cycle list|，需使用键

    |cycle list/BuGn-5|。

\noindent 这会初始化并选中相应的 |cycle list|。它也会初始化（但不选中）相应的
|colormap|。


注意，|colorbrewer| 自带的 |cycle list| 仅由\emph{颜色}组成。然而，一个好的
|cycle list| 通常还应附带标记，或许还有线型变化。要将纯颜色的 |cycle list| 与
标记组合，应使用 |cycle multi list|、|cycle multiindex list| 与
|cycle multiindex* list| 等特性，例如
%
\begin{codeexample}[code only]
\pgfplotsset{
    % initialize Set1-5:
    cycle list/Set1-5,
    % combine it with 'mark list*':
    cycle multiindex* list={
        mark list*\nextlist
        Set1-5\nextlist
    },
}
\end{codeexample}
%
\noindent |cycle multiindex* list| 的细节请参考参考手册。


\subsection{顺序方案}

顺序方案适用于从低到高有序的数据。

\noindent
\begin{tabular}{rrr}
    \MATRIXseq{BuGn}   & \MATRIXseq{PuRd}   & \MATRIXseq{Blues}   \\
    \MATRIXseq{BuPu}   & \MATRIXseq{RdPu}   & \MATRIXseq{Greens}  \\
    \MATRIXseq{GnBu}   & \MATRIXseq{YlGn}   & \MATRIXseq*{Greys}  \\
    \MATRIXseq{OrRd}   & \MATRIXseq{YlGnBu} & \MATRIXseq{Oranges} \\
    \MATRIXseq{PuBu}   & \MATRIXseq{YlOrBr} & \MATRIXseq{Purples} \\
    \MATRIXseq{PuBuGn} & \MATRIXseq{YlOrRd} & \MATRIXseq{Reds}    \\
\end{tabular}


\subsection{发散方案}

发散方案用于突出相对于某个均值的变化。

\noindent
\begin{tabular}{rrr}
    \MATRIXdiv{BrBG}   & \MATRIXdiv*{RdGy}   & \MATRIXdiv{RdYlBu} \\
    \MATRIXdiv{PiYG}   & \MATRIXdiv{PuOr}  & \MATRIXdiv{RdYlGn}\\
    \MATRIXdiv{PRGn}   & \MATRIXdiv{RdBu}  & \MATRIXdiv{Spectral}\\
\end{tabular}

注意，你可以调整 |point meta min| 与 |point meta max|，使 |colormap| 的均值与数据
相匹配（例如强制 |point meta min=-2| 与 |point meta max=+2|）。


\subsection{定性方案}

定性方案适用于颜色没有特殊顺序、但需要彼此可区分的情况。

\noindent
\begin{tabular}{rrr}
    \MATRIXqual{8}{Accent}    & \MATRIXqual{9}{Pastel1}[3] & \MATRIXqual{8}{Pastel2} \\
    \MATRIXqual{8}{Dark2}[1]  & \MATRIXqual{9}{Set1}       & \MATRIXqual{8}{Set2}[4]\\
    \MATRIXqual{12}{Paired}   & \MATRIXqual{12}{Set3}      & \\
\end{tabular}


\subsection{与 ColorBrewer 网站的交互}

要在 \url{http://colorbrewer2.org/} 上查找某个方案，例如上面选择的
\texttt{GnBu-7}，
%
\begin{itemize}
    \item 首先选择“Nature of your data”类型；本例为 “sequential”；
    \item 然后选择“Number of data classes”，即 7；
    \item 最后在“Pick a color scheme”部分选择相应方案。（这有点棘手，因为只能在
        右下角顶部看到名称，那里列出了所有颜色。选对之后应能看到“7-class GnBu”。
        不过，若知道“Gn”代表 green、“Bu”代表 blue，你就是在寻找一个从绿到蓝
        的方案，可能会有所帮助——本例中它是“Multi-hue”第一行的第三个配色方案。）
    \item 不过，颜色名（例如 \texttt{GnBu-H}）无法直接从页面提取，因为你不能简单
        地数到该数字并用对应字母替换。这是有意避免的，否则同一种颜色会以不同
        名称被定义多次，例如 \texttt{GnBu-H} 会被定义 3 次。
    \item 如果你需要方案 \texttt{GnBu-7} 的第 5 个颜色，不必反复试错或查阅手册。%
        \pgfplotsset{colormap/GnBu-7}%
        可以通过 \verb|index of colormap| 特性直接从相应色图中提取该颜色，例如

        |\tikz \fill[index of colormap={4 of GnBu-7}] (0,0) rectangle (1em,2ex);|

        其结果为 \tikz \fill[index of colormap={4 of GnBu-7}] (0,0)
        rectangle (1em,2ex);。注意索引为 $4$ 而非 $5$，因为索引从 $0$ 开始。
        另请注意，在使用该键之前必须先调用
        |\pgfplotsset{colormap/GnBu-7}|。
\end{itemize}


\subsection{外部示例}

若想查看一些 Brewer 本人使用其方案的示例，请访问
\url{https://www.census.gov/population/www/cen2000/atlas/index.html}。

\end{pgfplotslibrary}

\begin{tikzlibrary}{colorbrewer}
    本库仅包含形如 |GnBu-B| 的颜色定义。可用颜色列表见
    第~\ref{sec:pgfplots:brewer:usage} 节。
\end{tikzlibrary}
\endgroup
\endinput
